% chapters/13/praktikum.tex
\begin{praktikumbox}
Praktikum ini mengikuti Tugas Praktikum 10 pada RPS: membuat level sederhana
yang memuat lantai, rintangan, dan pemain yang bergerak serta bertabrakan
dengan lingkungannya.

\begin{langkah}
  \item \textbf{Susun levelnya secara \istilah{greyboxing}.} Buat lantai dari
        \istilah{Tilemap} atau dari balok, lalu tambahkan sekurang-kurangnya
        tiga bongkahan pada ketinggian yang berbeda dan dua dinding.

  \item \textbf{Tetapkan ukurannya lebih dahulu.} Tentukan tinggi lompatan
        pemain, lebar celah yang dapat dilewati, dan tinggi bongkahan. Susun
        levelnya agar seluruh celahnya terjangkau, dengan menyisakan kelonggaran
        sekitar sepuluh sampai dua puluh persen.

  \item \textbf{Tambahkan pemain.} Buat sebuah benda dengan
        \istilah{Rigidbody2D}, \istilah{Capsule Collider 2D}, dan skrip
        \code{KontrolPemain} dari Daftar~\ref{lst:ch13-pemain}. Atur
        \istilah{Constraints} agar putarannya terkunci.

  \item \textbf{Atur pemeriksaan tanahnya.} Buat sebuah benda kosong sebagai
        penanda di bawah kaki pemain, lalu hubungkan pada peubah
        \code{groundCheck}. Tentukan \istilah{layer} tanah, lalu pastikan
        lingkaran pemeriksaannya tampak pada panel \istilah{Scene} dan letaknya
        tepat.

  \item \textbf{Jalankan dan mainkan.} Pastikan pemain dapat berjalan,
        melompat, mendarat di atas bongkahan, dan terhalang oleh dinding.
        Apabila ada yang tidak sesuai, telusurilah menurut urutan pada
        Bagian~\ref{sec:masalah-fisika}, dan catat hasilnya.

  \item \textbf{Tambahkan daerah trigger.} Buat sebuah daerah pada ujung level,
        tandai \istilah{Collider}-nya sebagai \istilah{Is Trigger}, lalu
        tempelkan skrip \code{TriggerPintu} dari Daftar~\ref{lst:ch13-trigger}.
        Periksa pesan yang muncul pada panel \istilah{Console}.

  \item \textbf{Tambahkan kamera yang mengikuti.} Tulis skrip kamera seperti
        pada Bagian~\ref{sec:kamera-rasa}, lalu atur kelambatannya. Amati
        perbedaan rasa gerak pada dua nilai kelambatan yang berbeda.

  \item \textbf{Tulis catatan praktikum} yang memuat ukuran yang Anda tetapkan,
        tangkapan layar levelnya, catatan masalah fisika yang Anda temui beserta
        perbaikannya, dan nilai pengaturan yang Anda pilih.

  \item \textbf{Simpan pekerjaan} ke repositori Git dengan pesan penyimpanan
        yang menerangkan isinya.
\end{langkah}
\end{praktikumbox}

\begin{table}[htbp]
\centering
\caption{Lembar kerja Praktikum 10}
\label{tab:lembar-praktikum-10}
{\setlength{\tabcolsep}{4pt}\footnotesize%
\begin{tabularx}{\textwidth}{@{}L{4.4cm} Y@{}}
\toprule
\headrow \textbf{Butir yang dicatat} & \textbf{Keterangan} \\
\midrule
Tinggi lompatan, lebar celah, dan tinggi bongkahan & \\[1.2em]
\addlinespace
Masalah fisika yang ditemui beserta gejalanya & \\[2.4em]
\addlinespace
Penyebab dan perbaikan yang dilakukan & \\[2.4em]
\addlinespace
Nilai kelambatan kamera yang dipilih beserta alasannya & \\[1.2em]
\addlinespace
Pesan yang muncul pada panel \istilah{Console} ketika memasuki daerah trigger & \\[1.2em]
\bottomrule
\end{tabularx}}
\end{table}

\begin{catatabox}
Langkah kelima adalah langkah yang paling berharga bagi indikator \ind{c}.
Apabila level Anda berjalan mulus sejak percobaan pertama, tambahkan dengan
sengaja salah satu masalah dari Tabel~\ref{tab:masalah-fisika}---misalnya ubah
\istilah{Collision Detection} menjadi \istilah{Discrete}---lalu amati gejalanya
dan perbaiki kembali. Kebiasaan menimbulkan masalah dengan sengaja akan
memudahkan Anda mengenali masalah yang sama pada proyek akhir nanti.
\end{catatabox}
